% Respuesta a la pregunta 3 - Iñaki Cornejo - 822046480
\item ¿Por qué una empresa escoger una base de datos \textit{open source}?

  Las razones a continuación también aplican, en su mayor parte, a cualquier tipo de empresa que se ve presentada con la opción de elegir entre software propietario y software \textit{open source}. Una razón por la cual conviene elegir software de código abierto, es la autonomía y longevidad que este conlleva. Al usar software cerrado, una empresa se pone al merced los creadores de dicho software, si estos deciden dejar de desarrollarlo o revocar el acceso a este (claro a menos que esto vaya en contra al contrato de consumidor) entonces la empresa se ve obligada a lidiar con esto modificando su operación para hacer uso de otro software (lo cual puede ser extremadamente costoso y muchas veces resulta imposible). En cambio, si la organización encargada de desarrollar un software \textit{open source} deja de desarrollarlo es imposible que prohiban el uso del software. Esto va relacionado a la siguiente razón: flexibilidad. Si el software no se ajusta exactamente a las necesidades de una empresa, al tener disponibilidad del código resulta viable modificarlo para un uso específico. Por otro lado, si es software propietario no importa que tan pequeña sea la modificación, la única forma de conseguirla es comunicándoselo a los creadores del software y teniendo la suerte de ser escuchado. Por último, ya sea por adopción, por velocidad de desarrollo, por transparencia o por una de las miles razones posibles, todo el mejor software de administración de datos ¡es \textit{open source}! (sí Oracle, es la verdad)
